\section{Case Study: Serengeti Ecosystem Assessment}


We deployed \ConservAI{} across the full Snapshot Serengeti camera trap network (225 stations) for a 12-month evaluation period (January--December 2025).

\subsection{Key Findings}

\begin{enumerate}[leftmargin=*]
    \item \textbf{Population trends}: Wildebeest population estimated at 1.34M ($\pm$104K, 95\% CI), consistent with aerial survey estimates of 1.31M. Lion population showed a 4.2\% decline relative to the 2023 baseline.
    \item \textbf{Behavioral shifts}: TCA analysis revealed a 23-minute shift in peak activity time for Thomson's gazelle near tourist corridors, suggesting behavioral adaptation to human presence.
    \item \textbf{Habitat degradation}: The habitat quality index identified a 12\% decline in riparian corridor quality in the northern sector, correlating with upstream agricultural expansion detected in Sentinel-2 imagery.
    \item \textbf{Processing efficiency}: The full dataset of 1.2M images from the 12-month period was processed in 23.6 hours on a single V100 GPU, compared to an estimated 4,800 hours of manual expert classification.
\end{enumerate}

\subsection{Conservation Impact}

Based on \ConservAI{} findings, the Serengeti National Park Authority implemented three targeted interventions:

\begin{itemize}[leftmargin=*]
    \item Adjusted tourist vehicle routing to reduce behavioral disturbance in identified hotspots.
    \item Initiated a riparian restoration program in the degraded northern corridor.
    \item Increased anti-poaching patrols in areas where lion population decline was concentrated.
\end{itemize}

